% =============================================================
%  03_liste_ricorsione.tex
%  Capitolo 3 — Liste e ricorsione
% =============================================================

\chapter{Liste e ricorsione}

\begin{intuizione}
Le liste sono il cuore di Lisp. Il nome stesso viene da
\textit{LISt Processing}. Capire le liste, come sono fatte,
come si attraversano, come si costruiscono, è capire Lisp.
La ricorsione non è un'alternativa ai cicli: è il modo naturale
di pensare sulle liste.
\end{intuizione}

\vspace{4pt}

% ================================================================
\section{Cos'è una lista}
% ================================================================

\begin{concetto}{Una lista è una catena di cons-cell}
\keyline{}{Una lista non è un array. È una catena di coppie collegate.}

\detail{Ogni coppia si chiama \textbf{cons-cell} e contiene due slot:
\texttt{car} (il valore) e \texttt{cdr} (il puntatore alla prossima cella).
L'ultima cella punta a \texttt{nil}.}
\end{concetto}

\begin{esempio}[Struttura interna di una lista]
\begin{lstlisting}
;; La lista (a b c) in memoria:
;;
;;  [car|cdr] -> [car|cdr] -> [car|cdr] -> nil
;;    a            b            c
;;
;; Si costruisce con cons:
(cons 'a (cons 'b (cons 'c nil)))   ; => (A B C)

;; Equivalente:
(list 'a 'b 'c)   ; => (A B C)
\end{lstlisting}
\end{esempio}

\begin{confronto}
\textbf{Lista vs array:}
In Java \texttt{int[] a = \{1,2,3\}} alloca tre celle contigue in memoria.
In Lisp \texttt{(list 1 2 3)} crea tre cons-cell collegate.
Conseguenza: accesso al primo elemento O(1), accesso all'n-esimo O(n).
Le liste Lisp si usano per sequenze da attraversare dall'inizio, non per
accesso casuale.
\end{confronto}

\argomento{Le operazioni primitive: \fn{car}, \fn{cdr}, \fn{cons}}

\begin{concetto}{\fnt{car}, \fnt{cdr}, \fnt{cons}: le tre primitive}
\keyline{\fn{car}}{restituisce il \emph{primo elemento} (la testa).}
\keyline{\fn{cdr}}{restituisce il \emph{resto} (la coda, senza il primo).}
\keyline{\fn{cons}}{\emph{costruisce} una nuova cons-cell: \fn{(cons x lst)}.}

\boxsep
\detail{\fn{first}, \fn{rest} sono alias leggibili di \fn{car} e \fn{cdr}.
Usali nel codice nuovo: sono più chiari.}
\end{concetto}

\begin{esempio}[\fn{car}, \fn{cdr}, \fn{cons}]
\begin{lstlisting}
(car  '(a b c))   ; => A
(cdr  '(a b c))   ; => (B C)
(car  '(a))       ; => A
(cdr  '(a))       ; => NIL   (coda di una lista da un elemento)
(car  nil)        ; => NIL
(cdr  nil)        ; => NIL

;; cons aggiunge un elemento in testa
(cons 'a '(b c))   ; => (A B C)
(cons 'a nil)      ; => (A)
(cons 'a '())      ; => (A)   ('() = nil)

;; Alias leggibili
(first  '(a b c))   ; => A      (identico a car)
(rest   '(a b c))   ; => (B C)  (identico a cdr)
(second '(a b c))   ; => B
(third  '(a b c))   ; => C
\end{lstlisting}
\end{esempio}

\begin{attenzione}
\textbf{\fn{cons} non modifica la lista originale.}
\texttt{(cons 'x lst)} crea una \emph{nuova} cons-cell che punta a \texttt{lst}.
La lista \texttt{lst} rimane intatta.
Questo vale per tutte le operazioni funzionali sulle liste: non
modificano mai la struttura originale, restituiscono strutture nuove.
\end{attenzione}

\argomento{Composizione di \fn{car} e \fn{cdr}}

\begin{sintassi}
(cadr lst)    ; = (car (cdr lst))    = secondo elemento\\
(caddr lst)   ; = (car (cddr lst))   = terzo elemento\\
(cadddr lst)  ; = quarto elemento\\
(cddr lst)    ; = (cdr (cdr lst))    = lista senza i primi due
\end{sintassi}

\begin{esempio}[Accesso a elementi interni]
\begin{lstlisting}
(cadr   '(a b c d))   ; => B   (secondo)
(caddr  '(a b c d))   ; => C   (terzo)
(cadddr '(a b c d))   ; => D   (quarto)

;; In alternativa, piu' leggibile:
(nth 0 '(a b c))   ; => A
(nth 2 '(a b c))   ; => C
\end{lstlisting}
\end{esempio}

\begin{nota}
Le forme \fn{cadr}, \fn{caddr}, ecc. sono abbreviazioni storiche ancora
molto usate nel codice Lisp. \fn{nth} è più leggibile ma meno
idiomatico. In pratica: usa \fn{first}/\fn{second}/\fn{third} per le
prime posizioni, \fn{nth} per posizioni arbitrarie.
\end{nota}

\begin{workbook}
\textbf{Esercizi 22--26}
(Costruire liste, Accesso con car/cdr, Primo e resto,
Primo e ultimo, Scambia i primi due)\\[4pt]
Sono esercizi di manipolazione diretta: niente ricorsione ancora,
solo \fn{car}, \fn{cdr}, \fn{cons}, \fn{list}.
L'obiettivo è interiorizzare la struttura a cons-cell.
Falli nel REPL senza usare funzioni built-in come \fn{first}/\fn{last}/\fn{nth}.
\end{workbook}

% ================================================================
\section{Il pattern ricorsivo sulle liste}
% ================================================================

\begin{intuizione}
Ogni funzione che attraversa una lista segue lo stesso schema.
Impararlo una volta significa saperlo applicare a qualsiasi problema.
\end{intuizione}

\begin{concetto}{Lo schema base della ricorsione su liste}
\keyline{}{Caso base: lista vuota \fn{nil} → risultato terminale.}
\keyline{}{Caso ricorsivo: elabora \fn{(car lst)}, poi ricorri su \fn{(cdr lst)}.}

\boxsep
\detail{Ogni lista è o vuota (\fn{nil}) o una cons-cell con testa e coda.
La ricorsione rispecchia esattamente questa struttura.}

\boxsep
\detail{\textbf{Nota:} questi pattern assumono \emph{liste proprie}, catene
di cons-cell che terminano con \fn{nil}. Le liste improprie (dotted list)
come \fn{'(a b . c)} non terminano con \fn{nil} e richiedono gestione diversa.
In questo libro lavoriamo sempre con liste proprie.}
\end{concetto}

\begin{sintassi}
(defun f (lst)\\
\hspace{2em}(if (null lst)\\
\hspace{4em}... ; caso base: lista vuota\\
\hspace{4em}... ; caso ricorsivo: usa (car lst) e (f (cdr lst))
))
\end{sintassi}

\argomento{Tre pattern fondamentali}

Quasi tutte le funzioni sulle liste ricadono in uno di questi tre schemi:

\begin{esempio}[Pattern 1: calcola un valore (reduce)]
\begin{lstlisting}
;; Somma tutti gli elementi
(defun sum-list (lst)
  (if (null lst)
      0                              ; caso base: lista vuota = 0
      (+ (car lst)                   ; elabora la testa
         (sum-list (cdr lst)))))     ; ricorri sul resto

(sum-list '(1 2 3 4))   ; => 10

;; Stessa struttura: lunghezza
(defun length* (lst)
  (if (null lst)
      0
      (1+ (length* (cdr lst)))))

(length* '(a b c))   ; => 3
\end{lstlisting}
\end{esempio}

\begin{esempio}[Pattern 2: trasforma (map)]
\begin{lstlisting}
;; Aggiunge 1 a ogni elemento
(defun incrementa-tutti (lst)
  (if (null lst)
      nil                                    ; lista vuota -> lista vuota
      (cons (1+ (car lst))                   ; trasforma la testa
            (incrementa-tutti (cdr lst)))))  ; ricorri sul resto

(incrementa-tutti '(1 2 3))   ; => (2 3 4)

;; Stessa struttura: raddoppia ogni elemento
(defun raddoppia (lst)
  (if (null lst)
      nil
      (cons (* 2 (car lst))
            (raddoppia (cdr lst)))))
\end{lstlisting}
\end{esempio}

\begin{esempio}[Pattern 3: filtra (filter)]
\begin{lstlisting}
;; Mantieni solo i pari
(defun solo-pari (lst)
  (cond ((null lst) nil)             ; caso base
        ((evenp (car lst))           ; se la testa soddisfa il predicato:
         (cons (car lst)             ;   includila nel risultato
               (solo-pari (cdr lst))))
        (t                           ; altrimenti:
         (solo-pari (cdr lst)))))    ;   salta la testa

(solo-pari '(1 2 3 4 5 6))   ; => (2 4 6)
\end{lstlisting}
\end{esempio}

\begin{intuizione}
Questi tre pattern — \textbf{reduce}, \textbf{map}, \textbf{filter} —
sono il fondamento di tutta la programmazione funzionale.
Nel Capitolo~4 vedremo che \fn{reduce}, \fn{mapcar} e \fn{remove-if}
sono proprio questi pattern incapsulati in funzioni di ordine superiore.
Per ora, scriverli a mano con la ricorsione esplicita è fondamentale
per capire cosa fanno davvero.
\end{intuizione}

\begin{confronto}
\textbf{Da dove vengono i nomi map, reduce, filter?}

\medskip
\textbf{Map} viene dalla matematica: una funzione che \emph{mappa} ogni
elemento di un insieme in un elemento di un altro.
In Lisp il nome esiste dal 1960 (\fn{mapcar} in LISP 1.5).
Python lo adottò negli anni '90 con \texttt{map()};
JavaScript con \texttt{Array.map()} nel 2009.

\medskip
\textbf{Reduce} (detto anche \emph{fold} o \emph{accumulate}) viene
dall'idea di \emph{ridurre} una sequenza a un singolo valore applicando
un'operazione binaria cumulativamente.
In Common Lisp si chiama \fn{reduce}; in Haskell \texttt{foldl}/\texttt{foldr};
in Python \texttt{functools.reduce()}.
Il nome \emph{MapReduce} del framework di Google (2004) prende esattamente
questi due concetti per il calcolo distribuito su larga scala.

\medskip
\textbf{Filter} è il nome più ovvio: \emph{filtra} una sequenza tenendo
solo gli elementi che soddisfano un predicato.
In Common Lisp si chiama \fn{remove-if-not} o \fn{remove-if}
(per la variante che scarta); in Python \texttt{filter()};
in Java \texttt{Stream.filter()}.

\medskip
Questi tre pattern non sono un'invenzione di nessun linguaggio moderno.
Vengono da Lisp, che li ha avuti sin dall'inizio.
Quando vedi \texttt{.map().filter().reduce()} in JavaScript o
\texttt{stream().map().filter().collect()} in Java, stai usando
idee nate in Lisp sessant'anni prima.
\end{confronto}

\begin{workbook}
\textbf{Esercizi 27--38}
(Lunghezza, Somma, Prodotto, Massimo, Appartenenza, Conta occorrenze,
Posizione, Ultimo, Incrementa tutti, Filtra pari, Elimina tutte, Seleziona dispari)\\[4pt]
Ogni esercizio è uno dei tre pattern. Prima di scrivere il codice,
identifica quale: reduce, map o filter?\\[4pt]
Priorità agli essenziali: \textbf{27, 31, 35}.
\end{workbook}

% ================================================================
\section{Costruire liste: \fn{cons} vs \fn{append}}
% ================================================================

\begin{concetto}{\fnt{cons} e \fnt{append}: quando usare quale}
\keyline{\fn{cons}}{aggiunge \emph{un elemento} in testa. O(1).}
\keyline{\fn{append}}{concatena \emph{due liste}. O(n) sulla prima lista.}

\boxsep
\detail{Errore comune: usare \fn{append} quando basta \fn{cons},
o usare \fn{cons} con una lista quando serve \fn{append}.}
\end{concetto}

\begin{esempio}[\fn{cons} vs \fn{append}]
\begin{lstlisting}
;; cons: un elemento + una lista
(cons 'a '(b c))          ; => (A B C)   corretto
(cons '(a b) '(c d))      ; => ((A B) C D)  la testa e' una lista!

;; append: lista + lista
(append '(a b) '(c d))    ; => (A B C D)  corretto
(append '(a) '(b c))      ; => (A B C)

;; Errore tipico: cons dove serve append
(cons '(a b) '(c d))      ; => ((A B) C D)  NON e' quello che vuoi
(append '(a b) '(c d))    ; => (A B C D)    questo si'
\end{lstlisting}
\end{esempio}

\begin{esempio}[Implementare \fn{append} con ricorsione]
\begin{lstlisting}
;; append* ricorre su lst1, ricostruisce la catena
(defun append* (lst1 lst2)
  (if (null lst1)
      lst2                              ; caso base: lista1 vuota
      (cons (car lst1)                  ; testa di lst1
            (append* (cdr lst1) lst2))));; ricorri

(append* '(a b) '(c d))   ; => (A B C D)
\end{lstlisting}
\end{esempio}

\begin{nota}
\fn{append} è O(n) sulla lunghezza della \emph{prima} lista perché deve
ricopiare tutta la catena di cons-cell. La seconda lista non viene
copiata: viene solo puntata dall'ultima cella della prima.
Conseguenza pratica: in un loop che accumula risultati, è più efficiente
accumulare in testa con \fn{cons} e fare \fn{reverse} alla fine,
piuttosto che fare \fn{append} ad ogni passo.
\end{nota}

\begin{workbook}
\textbf{Esercizio 40} (Append)\\[4pt]
Implementa \fn{append*} da zero. È un esercizio $\bigstar\bigstar\bigstar$: il suggerimento
nel workbook aiuta a capire perché la ricorsione va su \fn{lst1} e non
su \fn{lst2}.
\end{workbook}

% ================================================================
\section{Reverse e l'accumulatore}
% ================================================================

\begin{concetto}{Il problema di \fnt{reverse} ricorsivo naif}
\keyline{}{La versione ovvia di reverse è $O(n^2)$ a causa di \fn{append}.}
\end{concetto}

\begin{esempio}[Reverse naif — $O(n^2)$]
\begin{lstlisting}
;; Versione ovvia: LENTA
(defun reverse-naif (lst)
  (if (null lst)
      nil
      (append (reverse-naif (cdr lst))   ; ricorri sul resto
              (list (car lst)))))         ; metti la testa in fondo

;; Perche' e' O(n^2)?
;; Al primo livello: append di una lista lunga n-1
;; Al secondo livello: append di una lista lunga n-2
;; ... -> n + (n-1) + ... + 1 = O(n^2)
\end{lstlisting}
\end{esempio}

\begin{concetto}{L'accumulatore: il pattern per O(n)}
\keyline{}{Porta un \emph{accumulatore} come parametro aggiuntivo.}
\keyline{}{Ad ogni passo, aggiungi in testa all'accumulatore.}
\keyline{}{A lista esaurita, l'accumulatore contiene il risultato.}

\boxsep
\detail{L'accumulatore parte vuoto e cresce man mano che la lista
originale si svuota. Alla fine, è già nell'ordine giusto senza
bisogno di \fn{append}.}
\end{concetto}

\begin{esempio}[Reverse con accumulatore — O(n)]
\begin{lstlisting}
(defun reverse* (lst)
  (labels ((rev-acc (lst acc)
             (if (null lst)
                 acc                 ; caso base
                 (rev-acc (cdr lst)
                          (cons (car lst) acc)))))
    (rev-acc lst nil)))              ; parte con acc vuoto

;; Traccia dell'esecuzione:
;; (reverse* '(a b c))
;; (rev-acc '(a b c) nil)
;; (rev-acc '(b c)   '(a))
;; (rev-acc '(c)     '(b a))
;; (rev-acc nil      '(c b a))   -> '(c b a)

(reverse* '(a b c d))   ; => (D C B A)
\end{lstlisting}
\end{esempio}

\begin{intuizione}
L'accumulatore è il modo funzionale di simulare una variabile locale
che accumula risultati in un ciclo imperativo. Invece di:
\begin{lstlisting}
;; Imperativo (concettuale):
;; result = []
;; for x in lst: result = [x] + result
;; return result
\end{lstlisting}
si porta \fn{result} come parametro della funzione ricorsiva.
\end{intuizione}

\begin{workbook}
\textbf{Esercizi 39, 51--53}\\[4pt]
Es.~39: scrivi \fn{reverse*} prima nella versione naif, poi con
l'accumulatore. Confronta le due: capire la differenza è fondamentale.\\[4pt]
Es.~51--53: fattoriale, reverse e somma-range in forma tail-recursive
con accumulatore. Sono esercizi su \fn{labels} --- il modo Lisp per
definire funzioni locali ricorsive.
\end{workbook}

% ================================================================
\section{Tail recursion}
% ================================================================

\begin{concetto}{Cos'è la tail recursion}
\keyline{}{Una chiamata ricorsiva è in \emph{posizione di coda} se è l'ultima\\
cosa che la funzione fa: nessun calcolo dopo di essa.}

\boxsep
\detail{Quando l'implementazione applica la \emph{tail call optimization} (TCO),
può riusare il frame dello stack invece di crearne uno nuovo.
Il risultato: ricorsione profonda con consumo di stack O(1).
Lo standard ANSI Common Lisp non garantisce TCO, ma le implementazioni
principali (SBCL, Allegro CL, CCL) la applicano sul codice compilato.}

\boxsep
\detail{$\triangleright$ \emph{Non sai cosa sono frame dello stack e stack overflow?
Vedi l'Appendice~\ref{app:stack} --- ``Lo stack delle chiamate'':
spiega come funzionano le chiamate di funzione a basso livello,
cosa è un frame di attivazione, quando si verifica lo stack overflow
e come la TCO lo evita.}}
\end{concetto}

\begin{esempio}[Tail recursion vs non-tail recursion]
\begin{lstlisting}
;; NON tail-recursive: c'e' una moltiplicazione DOPO la chiamata
(defun fattoriale-naif (n)
  (if (= n 0)
      1
      (* n (fattoriale-naif (- n 1)))))  ; * aspetta il risultato

;; Tail-recursive: la chiamata ricorsiva e' l'ULTIMA operazione
(defun fattoriale (n)
  (labels ((fact-acc (n acc)
             (if (= n 0)
                 acc                          ; ritorna direttamente
                 (fact-acc (- n 1) (* n acc)))))  ; niente dopo
    (fact-acc n 1)))

(fattoriale 100000)   ; con TCO attiva, niente stack overflow
\end{lstlisting}
\end{esempio}

\begin{confronto}
\textbf{In Common Lisp} la tail-call optimization (TCO) non è garantita
dallo standard ANSI. Le implementazioni principali (SBCL, Allegro CL, CCL)
la applicano sul codice compilato quando la chiamata è in posizione di coda.
Se la ricorsione è molto profonda e l'implementazione non applica TCO
(es.\ codice interpretato), usa \fn{loop} o \fn{do} come alternativa.
\end{confronto}

% ================================================================
\section{Ricorsione su strutture annidate}
% ================================================================

\begin{concetto}{Liste come alberi}
\keyline{}{Una lista può contenere altre liste: diventa un albero.}
\keyline{}{Per attraversarla tutta, la ricorsione deve scendere\\
anche \emph{dentro} gli elementi, non solo sul \fn{cdr}.}

\boxsep
\detail{La struttura \texttt{(a (b (c d) e) f)} è un albero con radice
implicita. Per visitarla tutta servono \emph{due} casi ricorsivi:
uno sulla testa (se è una lista, scendi dentro) e uno sulla coda.}
\end{concetto}

% ---- Figure: lista piatta e lista annidata come albero ----
\vspace{8pt}
\begin{center}

% -- 1. Lista piatta (a b c) come catena di cons-cell --
\begin{tikzpicture}[
  cell/.style={draw=csBlue!50, fill=csBlue!5, minimum width=1.8cm,
               minimum height=0.65cm, font=\small\ttfamily},
  ptr/.style={-{Stealth[length=5pt]}, csBlue!70, thick},
  nil/.style={draw=black!30, fill=black!5, minimum width=0.9cm,
              minimum height=0.65cm, font=\small\ttfamily\color{black!50}}
]
\node[cell] (c1) at (0,0)   {a};
\node[cell] (c2) at (2.4,0) {b};
\node[cell] (c3) at (4.8,0) {c};
\node[nil]  (n)  at (7.0,0) {nil};
\draw[ptr] (c1.east) -- (c2.west);
\draw[ptr] (c2.east) -- (c3.west);
\draw[ptr] (c3.east) -- (n.west);
\end{tikzpicture}

\smallskip
{\small\itshape\color{black!55}
\texttt{(a b c)} — lista piatta: catena lineare di cons-cell}

\bigskip

% -- 2. Lista annidata (a (b c) d) come albero --
\begin{tikzpicture}[
  level 1/.style={sibling distance=4.8cm, level distance=1.4cm},
  level 2/.style={sibling distance=2.8cm, level distance=1.4cm},
  level 3/.style={sibling distance=1.8cm, level distance=1.4cm},
  level 4/.style={sibling distance=1.5cm, level distance=1.3cm},
  edge from parent/.style={draw=csBlue!55, -{Stealth[length=4pt]}, thick},
  inode/.style={draw=csBlue!50, fill=csBlue!8, circle,
                minimum size=0.72cm, font=\small\ttfamily, inner sep=1pt},
  leaf/.style={draw=csGreen!55, fill=csGreen!8, rectangle, rounded corners=3pt,
               minimum width=0.7cm, minimum height=0.6cm,
               font=\small\bfseries\ttfamily\color{csGreen!50!black}},
  nilnode/.style={draw=black!25, fill=black!4, rectangle,
                  minimum width=0.6cm, minimum height=0.5cm,
                  font=\scriptsize\ttfamily\color{black!40}}
]
\node[inode] {cons}
  child { node[leaf] {a} }
  child { node[inode] {cons}
    child { node[inode] {cons}
      child { node[leaf] {b} }
      child { node[inode] {cons}
        child { node[leaf] {c} }
        child { node[nilnode] {nil} }
      }
    }
    child { node[inode] {cons}
      child { node[leaf] {d} }
      child { node[nilnode] {nil} }
    }
  };
\end{tikzpicture}

\smallskip
{\small\itshape\color{black!55}
\texttt{(a (b c) d)} --- lista annidata: albero binario di cons-cell}

\bigskip

% -- 3. Legenda --
\begin{tikzpicture}
\node[draw=csBlue!50, fill=csBlue!8, circle,
      minimum size=0.55cm, font=\scriptsize] (n1) at (0,0) {};
\node[font=\small\color{black!65}, right=5pt of n1] {cons-cell (nodo)};
\node[draw=csGreen!55, fill=csGreen!8, rectangle, rounded corners=2pt,
      minimum width=0.55cm, minimum height=0.5cm,
      font=\scriptsize\bfseries] (n2) at (4.0,0) {};
\node[font=\small\color{black!65}, right=5pt of n2] {atomo (foglia)};
\node[draw=black!25, fill=black!4, rectangle,
      minimum width=0.55cm, minimum height=0.45cm,
      font=\scriptsize\ttfamily\color{black!40}] (n3) at (8.0,0) {};
\node[font=\small\color{black!65}, right=5pt of n3] {\texttt{nil}};
\end{tikzpicture}

\smallskip
{\small\color{black!55}
Colonna sinistra (freccia \emph{car}) = scendi nella testa.
Colonna destra (freccia \emph{cdr}) = avanza nella coda.
La ricorsione su lista piatta percorre solo il \emph{cdr};
la ricorsione su albero percorre entrambi.}

\end{center}
\vspace{6pt}

\begin{sintassi}
(defun f (tree)\\
\hspace{2em}(cond ((null tree) ...)         \textit{; albero vuoto}\\
\hspace{2em}\hspace{6em}((atom tree) ...)   \textit{; foglia (atomo)}\\
\hspace{2em}\hspace{6em}(t                  \textit{; nodo interno (lista)}\\
\hspace{4em}... (f (car tree)) ...          \textit{; scendi nella testa}\\
\hspace{4em}... (f (cdr tree)) ...)))       \textit{; scendi nella coda}
\end{sintassi}

\begin{esempio}[\fn{flatten*}: appiattire strutture annidate]
\begin{lstlisting}
(defun flatten* (tree)
  (cond ((null tree) nil)              ; albero vuoto
        ((atom tree) (list tree))      ; foglia: avvolgila in una lista
        (t (append                     ; nodo: concatena
             (flatten* (car tree))     ;   risultato della testa
             (flatten* (cdr tree)))))) ;   risultato della coda

(flatten* '(a (b (c d) e) f))   ; => (A B C D E F)
(flatten* '((1 2) (3 (4 5))))   ; => (1 2 3 4 5)
\end{lstlisting}
\end{esempio}

\begin{esempio}[Contare atomi e sommare in profondità]
\begin{lstlisting}
;; Conta tutti gli atomi a qualsiasi profondita'
(defun count-atoms (tree)
  (cond ((null tree) 0)
        ((atom tree) 1)                          ; foglia = 1 atomo
        (t (+ (count-atoms (car tree))
              (count-atoms (cdr tree))))))

(count-atoms '(a (b c) ((d))))   ; => 4


;; Somma tutti i numeri a qualsiasi profondita'
(defun sum-tree (tree)
  (cond ((null tree)  0)
        ((numberp tree) tree)                    ; foglia numerica
        ((atom tree)    0)                       ; foglia non numerica
        (t (+ (sum-tree (car tree))
              (sum-tree (cdr tree))))))

(sum-tree '(1 (2 3) ((4 5))))   ; => 15
\end{lstlisting}
\end{esempio}

\begin{intuizione}
Il pattern su alberi è una generalizzazione di quello sulle liste.
Su una lista piatta ricorri solo sul \fn{cdr}.
Su un albero ricorri sia sul \fn{car} (per scendere) che sul \fn{cdr}
(per avanzare). È lo stesso schema dei parser, dei compilatori,
dei motori di ragionamento: tutto ciò che ha struttura ricorsiva si
attraversa con una funzione ricorsiva della stessa forma.
\end{intuizione}

\begin{workbook}
\textbf{Esercizi 55--60}
(Flatten, Profondità, Conta atomi, Somma profonda,
Ricerca annidata, Sostituzione profonda)\\[4pt]
Sono esercizi sull'attraversamento di alberi. Il primo (\fn{flatten*},
es.~55) è il più importante: capirlo significa capire il pattern.
Gli altri seguono la stessa struttura con piccole variazioni.\\[4pt]
Attenzione all'es.~55: il caso base è \emph{due} casi base
(\fn{null} e \fn{atom}), non uno solo come nelle liste piatte.
\end{workbook}

% ================================================================
\section{Funzioni utili sulle liste}
% ================================================================

\detail{Conoscerle evita di reinventare la ruota, e di capire cosa stiamo
reinventando negli esercizi.}

\begin{esempio}[Le funzioni built-in più usate]
\begin{lstlisting}
;; Accesso
(length '(a b c))           ; => 3
(nth 1 '(a b c))            ; => B  (0-based)
(last  '(a b c))            ; => (C)  restituisce una lista!
(butlast '(a b c))          ; => (A B)

;; Costruzione
(append '(a b) '(c d))      ; => (A B C D)
(reverse '(a b c))          ; => (C B A)

;; Ricerca
(member 'b '(a b c))        ; => (B C)  restituisce la coda!
(member 'x '(a b c))        ; => NIL
(position 'b '(a b c))      ; => 1

;; Selezione
(remove 'a '(a b a c))      ; => (B C)
(remove-duplicates '(a b a)) ; => (B A)  tiene l'ultima occorrenza
(subseq '(a b c d) 1 3)     ; => (B C)  come lst[1:3]

;; Predicati
(null  '())     ; => T
(consp '(a))    ; => T  (lista non vuota)
(listp '())     ; => T  (include nil)
\end{lstlisting}
\end{esempio}

\begin{attenzione}
\textbf{\fn{member} non restituisce T o NIL.}
Restituisce la \emph{coda} della lista a partire dall'elemento trovato,
oppure \fn{NIL} se non c'è. Questo significa che è comunque usabile
come booleano (\fn{nil} = falso, coda non-nil = vero), ma se ti aspetti
\fn{T} rimarrai sorpreso.
\begin{lstlisting}
(member 'b '(a b c))   ; => (B C)  non T!
(if (member 'b '(a b c)) "trovato" "no")  ; => "trovato"  ok lo stesso
\end{lstlisting}
\end{attenzione}

\begin{attenzione}
\textbf{\fn{remove-duplicates} tiene l'ultima occorrenza.}
Per default scarta le occorrenze \emph{precedenti} di ogni elemento duplicato.
\begin{lstlisting}
(remove-duplicates '(a b a))              ; => (B A)  la prima A sparisce
(remove-duplicates '(a b a) :from-end t)  ; => (A B)  tiene la prima
\end{lstlisting}
\end{attenzione}

% ================================================================
\section{Ricorsione con accumulatori avanzati}
% ================================================================

\argomento{Fibonacci efficiente: due accumulatori}

\begin{concetto}{Quando servono più accumulatori}
\keyline{}{Alcuni algoritmi richiedono tenere \emph{più stati} in parallelo.}
\detail{Fibonacci naif è $O(2^n)$ perché ricalcola gli stessi valori.
Con due accumulatori \fn{a = F(k)} e \fn{b = F(k+1)} si avanza di
un passo per chiamata: O(n) senza strutture dati ausiliarie.}
\end{concetto}

\begin{esempio}[Fibonacci O(n) con due accumulatori]
\begin{lstlisting}
(defun fibonacci-fast (n)
  (labels ((fib-step (n a b)
             ;; a = F(k), b = F(k+1)
             (if (= n 0)
                 a                          ; F(0) = a corrente
                 (fib-step (- n 1) b (+ a b))))) ; avanza: a<-b, b<-a+b
    (fib-step n 0 1)))                      ; parte con F(0)=0, F(1)=1

(fibonacci-fast 10)   ; => 55
(fibonacci-fast 50)   ; => 12586269025   (nessun problema di stack)
\end{lstlisting}
\end{esempio}

\argomento{Il crivello di Eratostene: ricorsione + array}

% --- Spiegazione degli array prima dell'esempio ---

\begin{concetto}{Array in Common Lisp}
\keyline{\fn{make-array}}{crea un array di dimensione fissa.}
\keyline{\fn{aref}}{accede (legge o scrive) a un elemento per indice.}
\keyline{\fn{setf}\ \fn{aref}}{modifica un elemento in posizione.}

\boxsep
\detail{A differenza delle liste, un array alloca le celle in
memoria contigua. Accesso e modifica sono $O(1)$, qualunque sia
l'indice.}
\end{concetto}

\begin{tcolorbox}[
  enhanced,
  colback=black!2, colframe=black!20,
  arc=3pt, boxrule=0.4pt,
  left=12pt, right=12pt, top=9pt, bottom=9pt,
  before skip=8pt, after skip=10pt
]
\small\color{black!65}

\textbf{Cosa è un array in Common Lisp.}

\medskip
Un array è una sequenza di celle numerate a partire da 0,
allocate in memoria contigua.
La dimensione viene fissata alla creazione e non cambia.
A differenza della lista --- che è una catena di cons-cell distribuite
in memoria --- l'array permette di raggiungere qualsiasi elemento
direttamente, in tempo costante $O(1)$, conoscendo solo l'indice.

\medskip
\textbf{Creare un array.}
La funzione è \fn{make-array}. Il primo argomento è la dimensione;
\fn{:initial-element} imposta il valore iniziale di ogni cella.

\begin{tcolorbox}[
  enhanced, colback=clBg, colframe=gray!40,
  arc=2pt, boxrule=0.4pt,
  left=8pt, right=6pt, top=4pt, bottom=4pt,
  before skip=4pt, after skip=4pt
]
\begin{lstlisting}
;; Array di 5 elementi, tutti inizializzati a 0
(make-array 5 :initial-element 0)
; => #(0 0 0 0 0)

;; Array di 4 booleani, tutti T
(make-array 4 :initial-element t)
; => #(T T T T)

;; Leggerlo: la notazione #(...) e' la rappresentazione di un array
\end{lstlisting}
\end{tcolorbox}

\textbf{Leggere un elemento: \fn{aref}.}
\fn{(aref arr i)} restituisce l'elemento all'indice \fn{i} (0-based).

\begin{tcolorbox}[
  enhanced, colback=clBg, colframe=gray!40,
  arc=2pt, boxrule=0.4pt,
  left=8pt, right=6pt, top=4pt, bottom=4pt,
  before skip=4pt, after skip=4pt
]
\begin{lstlisting}
(defparameter *arr* (make-array 5 :initial-element 0))

(aref *arr* 0)   ; => 0   (primo elemento)
(aref *arr* 4)   ; => 0   (ultimo elemento)
(aref *arr* 5)   ; ERRORE: indice fuori range
\end{lstlisting}
\end{tcolorbox}

\textbf{Modificare un elemento: \fn{setf} + \fn{aref}.}
\fn{setf} è l'operatore di assegnazione generalizzato.
Combinato con \fn{aref} modifica una cella specifica.

\begin{tcolorbox}[
  enhanced, colback=clBg, colframe=gray!40,
  arc=2pt, boxrule=0.4pt,
  left=8pt, right=6pt, top=4pt, bottom=4pt,
  before skip=4pt, after skip=4pt
]
\begin{lstlisting}
(setf (aref *arr* 2) 42)   ; scrive 42 in posizione 2
(aref *arr* 2)             ; => 42
;; *arr* e' ora #(0 0 42 0 0)

;; Modificare piu' celle in un colpo solo
(setf (aref *arr* 0) 10
      (aref *arr* 1) 20)
;; *arr* e' ora #(10 20 42 0 0)
\end{lstlisting}
\end{tcolorbox}

\textbf{Differenza fondamentale con la lista.}
L'operazione \fn{(aref arr i)} è identica per $i = 0$ e $i = 1000$:
raggiunge la cella direttamente senza scorrere nulla.
Con una lista, \fn{(nth 1000 lst)} deve scorrere 1000 cons-cell.
Questa differenza è irrilevante per liste di 5 elementi;
diventa decisiva su strutture di migliaia o milioni di elementi.

\medskip
\textbf{Quando usare array invece di liste.}
\begin{itemize}
  \item Quando l'algoritmo richiede accesso per indice
        (``vai alla posizione $k$'')
  \item Quando la dimensione è nota a priori e non cambia
  \item Quando la struttura viene modificata in-place
        (algoritmi di marcatura, ordinamento, crivelli)
\end{itemize}
Le liste rimangono la scelta migliore per strutture di dimensione
variabile, costruite incrementalmente con \fn{cons}.
\end{tcolorbox}

\begin{concetto}{Quando uscire dalla ricorsione pura}
\keyline{}{Alcuni algoritmi richiedono \emph{accesso casuale} a posizioni\\
arbitrarie: per questi, un array è la struttura giusta.}
\keyline{}{L'uscita dalla ricorsione pura è una scelta \emph{locale}:\\
l'interfaccia esterna rimane funzionale.}

\boxsep
\detail{Il crivello di Eratostene marca celle di un array booleano.
La struttura interna è imperativa (array mutabile + loop),
ma l'interfaccia esterna è funzionale (restituisce una lista di primi).
Ricorsione e stato mutabile coesistono in Common Lisp.}
\end{concetto}

\begin{tcolorbox}[
  enhanced,
  colback=black!2, colframe=black!20,
  arc=3pt, boxrule=0.4pt,
  left=12pt, right=12pt, top=9pt, bottom=9pt,
  before skip=8pt, after skip=10pt
]
\small\color{black!65}

\textbf{Perché uscire dalla ricorsione pura?}

\medskip
La ricorsione è naturale quando il problema ha struttura ricorsiva:
una lista si scompone in testa e coda, un albero in radice e sottoalberi.
Ma non tutti i problemi hanno questa forma.

\medskip
Il crivello di Eratostene ne è l'esempio classico.
L'algoritmo funziona così: parti con un vettore di N celle, tutte
segnate ``primo''. Per ogni numero $i$ non ancora eliminato, cancella
tutti i suoi multipli $2i, 3i, 4i, \ldots$ Dal vettore che rimane,
raccogli i numeri non cancellati.

\medskip
Il punto critico è \textbf{cancellare un multiplo}: significa accedere
alla cella in posizione $j = k \cdot i$ e modificarla.
Con una lista Lisp questo costo è $O(j)$ per ogni accesso
(devi scorrere $j$ cons-cell). Con un array è $O(1)$: accesso diretto.
Su N numeri, la differenza tra $O(1)$ e $O(N)$ per ogni accesso
è enorme: l'algoritmo passerebbe da $O(N \log \log N)$ a $O(N^2)$.

\medskip
\textbf{La regola generale:}
quando l'algoritmo richiede di \emph{tornare indietro} o
\emph{saltare} a posizioni arbitrarie nella struttura dati, e non
di attraversarla dall'inizio alla fine, la lista è la struttura sbagliata.
In questi casi usa un array con \fn{make-array} e \fn{aref}.

\medskip
\textbf{Lo stile ibrido} è la risposta professionale:
usa la struttura mutabile \emph{localmente}, all'interno della funzione,
per fare il lavoro efficiente; esponi all'esterno solo il risultato
come lista o valore funzionale.
Il chiamante non sa e non deve sapere cosa c'è dentro.
Questo è incapsulamento: la scelta implementativa rimane privata.

\medskip
\textbf{In breve}: non è un fallimento del paradigma funzionale.
È una scelta pragmatica che Common Lisp rende naturale,
a differenza dei linguaggi funzionali puri (Haskell) che richiedono
strutture dati persistenti o monadi per gestire lo stato.
\end{tcolorbox}

\begin{esempio}[Crivello di Eratostene]
\begin{lstlisting}
(defun crivello (n)
  (let ((arr (make-array (1+ n) :initial-element t)))
    (setf (aref arr 0) nil
          (aref arr 1) nil)            ; 0 e 1 non sono primi
    (loop for i from 2 to (isqrt n) do ; per ogni potenziale primo
          (when (aref arr i)           ; se non e' gia' segnato
            (loop for j from (* i i) to n by i do  ; segna multipli
                  (setf (aref arr j) nil))))
    (loop for i from 2 to n           ; raccogli i rimasti
          when (aref arr i)
          collect i)))

(crivello 20)   ; => (2 3 5 7 11 13 17 19)
\end{lstlisting}
\end{esempio}

\begin{nota}
Questo è un esempio di \textbf{stile ibrido}: struttura dati mutabile
(array) usata localmente per efficienza, interfaccia esterna funzionale.
Il chiamante non sa (e non deve sapere) che dentro c'è un array.
È un pattern professionale molto comune in Lisp.
\end{nota}

\begin{workbook}
\textbf{Esercizi 54, 64, 65}
(Fibonacci fast, Crivello di Eratostene, Fattori primi)\\[4pt]
Questi tre esercizi richiedono ragionamento algoritmico, non solo
pattern ricorsivi. Consulta il suggerimento nel workbook
prima di tentare la soluzione.\\[4pt]
L'es.~64 (crivello) è il primo esercizio che combina \fn{loop},
\fn{make-array} e \fn{setf}: è un ottimo punto di contatto tra
stile funzionale e stile imperativo.
\end{workbook}

% ================================================================
\section*{Riepilogo del capitolo}
% ================================================================

\begin{tcolorbox}[
  enhanced,
  colback=csBlue!5, colframe=csBlue!40,
  arc=4pt, boxrule=0.6pt,
  left=14pt, right=14pt, top=12pt, bottom=12pt,
  before skip=16pt, after skip=16pt
]
\textbf{Quello che hai imparato:}

\medskip
\begin{itemize}
  \item \textbf{Struttura}: una lista è una catena di cons-cell.
        \fn{car} = testa, \fn{cdr} = coda, \fn{cons} = costruisce.
  \item \textbf{Pattern ricorsivi}: reduce (calcola un valore),
        map (trasforma), filter (seleziona).
        Sono i fondamenti della programmazione funzionale.
  \item \textbf{\fn{cons} vs \fn{append}}: \fn{cons} aggiunge in testa O(1),
        \fn{append} concatena O(n). Usare \fn{cons} + \fn{reverse}
        è più efficiente di \fn{append} in un loop.
  \item \textbf{Accumulatore}: porta il risultato parziale come parametro.
        Rende la ricorsione tail-recursive e O(n).
  \item \textbf{Tail recursion}: chiamata in posizione di coda,
        ottimizzabile dallo stack. Usa \fn{labels} per la funzione ausiliaria.
  \item \textbf{Alberi}: liste annidate richiedono due casi ricorsivi:
        sulla testa (scendi) e sulla coda (avanza).
  \item \textbf{Stile ibrido}: array per accesso casuale, lista come
        interfaccia esterna. Ricorsione e mutabilità coesistono.
\end{itemize}

\medskip
\textbf{Prossimo:} Capitolo~4, \emph{Higher-order e closure}.
Le funzioni come valori, \fn{mapcar}, \fn{reduce}, lambda, closure.
Quello che hai scritto a mano in questo capitolo diventa
un'unica chiamata.
\end{tcolorbox}

\begin{workbook}
\textbf{Prima di andare al Capitolo~4:}\\[4pt]
Completa gli esercizi \textbf{22--65} (capitoli \emph{Liste} e
\emph{Ricorsione avanzata}).\\[4pt]
Priorità assoluta: \textbf{27, 31, 35, 39, 40, 51, 55}.
Sono gli esercizi che insegnano i pattern: gli altri li applicano.\\[4pt]
Se hai difficoltà con un esercizio $\bigstar\bigstar\bigstar$ o superiore, leggi il
suggerimento prima della soluzione.
\end{workbook}
